\makeatletter
\@ifundefined{keepfig}{\newenvironment{keepfig}{\par\edef\keepfig@pi{\the\parindent}\noindent
  \begin{minipage}{\linewidth}\parindent=\keepfig@pi\relax}{\end{minipage}\par}}{}
\makeatother
\section{Вычисление групп Галуа над \texorpdfstring{$\Q$}{Q}}

В этом разделе описывается практический способ вычислять группы Галуа многочленов над $\Q$. Его основа~--- теорема
Дедекинда \rf{4.28}: разложение многочлена $f\in\Z[X]$ на неприводимые множители по модулю простого числа $p$ показывает
цикловой тип некоторого элемента $G_f$. Раскладывая $f$ по модулю многих простых чисел, мы находим много элементов $G_f$
и часто тем самым определяем всю группу.

Напомним обозначения: для сепарабельного $f\in F[X]$ через $F_f$ обозначается поле разложения $f$ над $F$, а через
$G_f=\Gal(F_f/F)$~--- группа Галуа $f$ (\rf{3.27a}). Группа $G_f$ переставляет корни $\alpha_1,\dots,\alpha_m$ многочлена $f$
($m=\deg f$), и это задаёт вложение $G_f\hookrightarrow\Sym\{\alpha_1,\dots,\alpha_m\}\cong S_m$.

\subsection*{Орбиты и неприводимые множители}

Следующее предложение обобщает \rf{4.5}: не только транзитивность, но и все орбиты группы $G_f$ видны по разложению $f$. По
существу это \rf{3.27d}; здесь оно доказывается иначе~--- через делители $f$, без продолжения изоморфизмов.

\begin{proposition}[4.26]
Пусть $f\in F[X]$~--- сепарабельный многочлен, и орбиты группы $G_f$ на множестве корней $f$ состоят из $m_1,\dots,m_r$
элементов. Тогда $f=f_1\cdots f_r$, где $f_i\in F[X]$ неприводимы и $\deg f_i=m_i$.
\end{proposition}

\begin{proof}
Можно считать $f$ унитарным. Пусть $R=\{\alpha_1,\dots,\alpha_m\}$~--- множество корней $f$ в $F_f$; так как $f$ сепарабелен,
$f=\prod_{\alpha\in R}(X-\alpha)$. По однозначности разложения в $F_f[X]$ унитарные делители $f$ в $F_f[X]$ находятся во
взаимно однозначном соответствии с подмножествами $S\subset R$:
\[
S\ \longleftrightarrow\ f_S=\prod_{\alpha\in S}(X-\alpha).
\]
Для $\sigma\in G_f$ (действующего на коэффициенты) $\sigma f_S=\prod_{\alpha\in S}(X-\sigma\alpha)=f_{\sigma S}$. Поэтому
$\sigma f_S=f_S$ при всех $\sigma\in G_f$ тогда и только тогда, когда $\sigma S=S$ при всех $\sigma$, т.\,е.\ $S$
устойчиво относительно $G_f$. С другой стороны, коэффициенты $f_S$ неподвижны относительно $G_f$ тогда и только тогда,
когда они лежат в $F$: поле $F_f^{G_f}$ равно $F$, так как $F_f/F$~--- расширение Галуа (\rf{3.10}). Итак, унитарные
делители $f$ в $F[X]$~--- это $f_S$ для $G_f$-устойчивых $S$. Неприводимые унитарные делители соответствуют минимальным
непустым устойчивым подмножествам, а это в точности орбиты: орбита устойчива, а всякое непустое устойчивое
подмножество есть объединение орбит. Разбиение $R$ на орбиты $O_1,\dots,O_r$ даёт $f=f_{O_1}\cdots f_{O_r}$ с
$\deg f_{O_i}=|O_i|=m_i$.
\end{proof}

\begin{remark}[4.27]
Доказательство даёт более точное утверждение: если $R=O_1\sqcup\dots\sqcup O_r$~--- разбиение множества корней на орбиты
$G_f$, то
\[
f=f_1\cdots f_r,\qquad f_i=\prod_{\alpha\in O_i}(X-\alpha),
\]
~--- разложение унитарного $f$ в произведение унитарных неприводимых многочленов из $F[X]$. В частности, $f$ неприводим
тогда и только тогда, когда $G_f$ действует на корнях транзитивно (\rf{3.27d}, \rf{4.5}).
\end{remark}

Применим это к конечному полю (цикловой тип перестановки и его связь с орбитами~--- см.\ абзац перед \rf{4.25a}).

\begin{remark}[4.27a]
Если $F=\F_q$ и $f=f_1\cdots f_r\in F[X]$~--- сепарабельный многочлен с неприводимыми множителями $f_i$, $\deg f_i=m_i$,
то $G_f$ порождена автоморфизмом Фробениуса $\varphi_q\colon x\mapsto x^q$; поэтому циклы $\varphi_q$ на корнях $f$~---
это орбиты $G_f$, т.\,е.\ по \rf{4.27} множества корней множителей $f_i$: \emph{цикловой тип $\varphi_q$ равен
$(m_1,\dots,m_r)$}. Это предложение \rf{4.25a} (там же пример).
\end{remark}

\subsection*{Теорема Дедекинда}

Пусть теперь $f\in\Z[X]$ унитарен. Его редукция $\bar f\in\F_p[X]$ по модулю $p$ имеет циклическую группу Галуа, и по \rf{4.25a}
её образующая~--- Фробениус~--- имеет цикловой тип, равный набору степеней неприводимых множителей $\bar f$. Теорема
Дедекинда утверждает, что тот же цикловой тип имеет некоторый элемент группы $G_f$ над $\Q$. Идея доказательства:
построить кольцо $A$, содержащее корни $f$, и его гомоморфизм на поле разложения $\bar f$, переводящий корни $f$ в корни
$\bar f$; затем показать, что каждый автоморфизм поля разложения $\bar f$ поднимается до элемента $G_f$, переставляющего
корни <<так же>>. Доказательство разбито на несколько шагов (\rf{4.27b}--\rf{4.27e}).

\emph{Обозначения до конца подраздела.} $f\in\Z[X]$~--- унитарный многочлен степени $m$, $E=\Q_f$~--- его поле
разложения, $\alpha_1,\dots,\alpha_m$~--- корни $f$ в $E$ (с повторениями, если они есть), $G=\Gal(E/\Q)$ и
\[
A=\Z[\alpha_1,\dots,\alpha_m]\subset E
\]
~--- наименьшее подкольцо $E$, содержащее все корни, т.\,е.\ множество значений $P(\alpha_1,\dots,\alpha_m)$ многочленов
$P\in\Z[X_1,\dots,X_m]$. Черта означает редукцию: $\bar f\in\F_p[X]$~--- многочлен с коэффициентами, приведёнными по
модулю $p$.

\begin{lemma}[4.27b]
\begin{enumerate}[label=\textup{(\alph*)}]
\item $A$~--- конечно порождённая абелева группа без кручения, т.\,е.\ $A\cong\Z^s$ для некоторого $s\ge1$; каждый
$\sigma\in G$ переводит $A$ в себя.
\item Никакое простое число $\ell$ не обратимо в $A$, и $A\cap\Q=\Z$.
\item Для каждого простого $p$ существует максимальный идеал $P$ кольца $A$, содержащий $p$. Для любого такого $P$
имеем $P\cap\Z=p\Z$, а $A/P$~--- конечное поле характеристики $p$, содержащее $\F_p=\Z/p\Z$ (образ $\Z$).
\end{enumerate}
\end{lemma}

\begin{proof}
(a) Так как $f(\alpha_i)=0$ и $f$ унитарен степени $m$, то $\alpha_i^m$~--- целочисленная линейная комбинация
$1,\alpha_i,\dots,\alpha_i^{m-1}$, а по индукции и каждая степень $\alpha_i^k$. Поэтому любой одночлен
$\alpha_1^{k_1}\cdots\alpha_m^{k_m}$ есть целочисленная комбинация одночленов с показателями $k_j<m$, и $A$ порождается как
абелева группа конечным множеством таких одночленов. Кручения нет, так как $A\subset E$, а $\operatorname{char}E=0$. По
теореме о строении конечно порождённых абелевых групп $A\cong\Z^s$, и $s\ge1$, так как $1\in A$. Элемент $\sigma\in G$
переставляет корни, поэтому $\sigma(P(\alpha_1,\dots,\alpha_m))=P(\sigma\alpha_1,\dots,\sigma\alpha_m)\in A$.

(b) Из $A\cong\Z^s$ следует $A/\ell A\cong(\Z/\ell\Z)^s\ne0$, т.\,е.\ $\ell A\ne A$. Если бы $\ell c=1$ для некоторого
$c\in A$, то $a=\ell(ca)\in\ell A$ для всех $a\in A$, т.\,е.\ $\ell A=A$. (Другое объяснение: из $1/\ell\in A$ следовало
бы $\Z[1/\ell]\subset A$, но $\Z[1/\ell]$ не является конечно порождённой абелевой группой, а подгруппы конечно порождённых
абелевых групп конечно порождены.) Пусть теперь $u/v\in A$, где $u,v\in\Z$ взаимно просты и $v>1$; пусть $\ell$~---
простой делитель $v$. Найдём $x,y\in\Z$ с $xu+yv=1$; тогда $1/v=x\cdot(u/v)+y\in A$ и $1/\ell=(v/\ell)\cdot(1/v)\in A$~---
противоречие. Значит, $A\cap\Q=\Z$.

(c) Кольцо $A/pA\cong(\Z/p\Z)^s$ конечно и ненулевое. Среди его собственных идеалов (например, нулевого) выберем
идеал наибольшего порядка; он максимален. Его прообраз $P$ в $A$~--- максимальный идеал, содержащий $p$ (идеалы $A/pA$
соответствуют идеалам $A$, содержащим $pA$). Пересечение $P\cap\Z$~--- идеал $\Z$, содержащий $p$ и не содержащий $1$; так
как $p\Z$~--- максимальный идеал $\Z$, то $P\cap\Z=p\Z$. Ядро композиции $\Z\to A\to A/P$ равно $P\cap\Z=p\Z$, так что
образ $\Z$ в поле $A/P$ есть $\F_p$. Наконец, $A/P$~--- фактор конечного кольца $A/pA$, т.\,е.\ конечно.
\end{proof}

Зафиксируем простое $p$ и максимальный идеал $P\supset p$ кольца $A$, и обозначим $a\mapsto\bar a$ отображение
$A\to A/P$. По \rf{4.27b}(c) $\F_p\subset A/P$, и редукция коэффициентов по модулю $p$~--- то же, что их редукция по
модулю $P$.

\begin{lemma}[4.27c]
\begin{enumerate}[label=\textup{(\alph*)}]
\item В $(A/P)[X]$ выполнено $\bar f=(X-\bar\alpha_1)\cdots(X-\bar\alpha_m)$, и $A/P=\F_p[\bar\alpha_1,\dots,\bar\alpha_m]$;
значит, $A/P$~--- поле разложения $\bar f$ над $\F_p$.
\item Дискриминант $D(f)=\prod_{i<j}(\alpha_i-\alpha_j)^2$ лежит в $\Z$, и $D(\bar f)=\overline{D(f)}$ (т.\,е.\
$D(f)\bmod p$).
\item Следующие условия равносильны: \textup{(i)} $p\nmid D(f)$; \textup{(ii)} $\bar f$ сепарабелен; \textup{(iii)}
элементы $\bar\alpha_1,\dots,\bar\alpha_m$ попарно различны. При этих условиях $f$ сепарабелен, а редукция
$\alpha_i\mapsto\bar\alpha_i$~--- биекция множества корней $f$ на множество корней $\bar f$.
\end{enumerate}
\end{lemma}

\begin{proof}
(a) Равенство $f=\prod_i(X-\alpha_i)$ выполнено в $E[X]$, а все его коэффициенты лежат в $A$, так что оно выполнено в
$A[X]$. Гомоморфизм $A\to A/P$ продолжается до гомоморфизма колец $A[X]\to(A/P)[X]$ (применение к коэффициентам);
он переводит $f$ в $\bar f$, а $\prod(X-\alpha_i)$~--- в $\prod(X-\bar\alpha_i)$. Каждый элемент $A$ имеет вид
$\sum k_{i_1\dots i_m}\alpha_1^{i_1}\cdots\alpha_m^{i_m}$ с $k_{i_1\dots i_m}\in\Z$, и его образ равен
$\sum\bar k_{i_1\dots i_m}\bar\alpha_1^{i_1}\cdots\bar\alpha_m^{i_m}$; поэтому $A/P$ порождается над $\F_p$ элементами
$\bar\alpha_i$, т.\,е.\ корнями $\bar f$. Это и есть определение поля разложения.

(b) $D(f)\in A$, и $D(f)$ неподвижен относительно $G$ (элементы $G$ переставляют корни, а $D(f)$~--- симметрическая функция
корней), значит, $D(f)\in E^G\cap A=\Q\cap A=\Z$ по \rf{4.27b}(b). Применив к $D(f)=\prod_{i<j}(\alpha_i-\alpha_j)^2$
гомоморфизм $A\to A/P$ и учитывая (a), получаем $\overline{D(f)}=\prod_{i<j}(\bar\alpha_i-\bar\alpha_j)^2=D(\bar f)$.

(c) Из (a) корни $\bar f$ в поле $A/P$~--- это $\bar\alpha_1,\dots,\bar\alpha_m$ с кратностями; поэтому (ii)$\iff$(iii), а
(iii)$\iff\prod_{i<j}(\bar\alpha_i-\bar\alpha_j)^2\ne0\iff$(i) по (b). Если $\bar\alpha_i$ различны, то различны и $\alpha_i$,
так что $f$ сепарабелен.
\end{proof}

Условие (iii) необходимо для всего дальнейшего: без него редукция склеивает корни, и связать перестановки корней $f$ и
$\bar f$ нельзя. Например, $\bar f=(X+1)^4$ для $f=X^4+1$ и $p=2$ (здесь $D(f)=256$), и разложение по модулю $2$ ничего
не говорит о $G_f$.

\begin{definition}[4.27d]
\emph{Группа разложения} идеала $P$~--- подгруппа
\[
D_P=\{\sigma\in G\mid\sigma(P)=P\}.
\]
Каждый $\sigma\in D_P$ переводит $A$ в $A$ и $P$ в $P$, поэтому корректно определён автоморфизм
$\bar\sigma\colon A/P\to A/P$, $\bar\sigma(\bar a)=\overline{\sigma(a)}$ (если $a-a'\in P$, то $\sigma a-\sigma a'\in\sigma P=P$);
он тождествен на $\F_p$ и потому лежит в $G_{\bar f}=\Gal\bigl((A/P)/\F_p\bigr)$.
\end{definition}

Ясно, что $D_P$~--- подгруппа, а $\theta\colon D_P\to G_{\bar f}$, $\sigma\mapsto\bar\sigma$,~--- гомоморфизм групп:
$\overline{\sigma\tau}(\bar a)=\overline{\sigma(\tau a)}=\bar\sigma(\bar\tau(\bar a))$. Заметим также, что для
$\sigma\in G$ образ $\sigma(P)$ и прообраз $\sigma^{-1}(P)$~--- снова максимальные идеалы $A$ (так как $\sigma|_A$~---
автоморфизм кольца $A$), и $\sigma^{-1}(P)=P\iff\sigma(P)=P\iff\sigma\in D_P$.

\begin{proposition}[4.27e]
Пусть $p\nmid D(f)$. Тогда:
\begin{enumerate}[label=\textup{(\alph*)}]
\item гомоморфизм $\theta\colon D_P\to G_{\bar f}$ является изоморфизмом;
\item существует единственный элемент $\sigma_P\in G$, для которого
\begin{equation}\label{eq:s6-frob}
\sigma_P(a)\equiv a^p\pmod P\qquad\text{для всех }a\in A;
\end{equation}
он лежит в $D_P$ и порождает $D_P$. Он называется \emph{элементом Фробениуса} идеала $P$.
\end{enumerate}
\end{proposition}

\begin{proof}
\emph{Инъективность $\theta$.} Если $\bar\sigma=\id$, то $\overline{\sigma\alpha_i}=\bar\sigma(\bar\alpha_i)=\bar\alpha_i$ для
всех $i$. Но $\sigma\alpha_i$~--- один из корней $\alpha_j$, а корни $f$ попарно несравнимы по модулю $P$ (\rf{4.27c}(c)); значит,
$\sigma\alpha_i=\alpha_i$ для всех $i$, и $\sigma=\id$, так как $E=\Q[\alpha_1,\dots,\alpha_m]$.

\emph{Сюръективность $\theta$.} Нужно показать, что каждый автоморфизм поля $A/P$ над $\F_p$ имеет вид $\bar\sigma$ с
$\sigma\in D_P$. Поле $A/P$ конечно, и его мультипликативная группа циклична; пусть $\bar a$~--- её образующая. Тогда
$A/P=\F_p[\bar a]$, $\bar a\ne0$. Выберем прообраз $a\in A$; так как $\bar a\ne0$, то $a\notin P$.

\emph{Шаг 1: элемент $b$, <<невидимый>> вне $D_P$.} Пусть $Q_1,\dots,Q_k$~--- различные идеалы среди $\sigma^{-1}(P)$,
$\sigma\in G\smallsetminus D_P$. Это максимальные идеалы, отличные от $P$, поэтому $P+Q_j=A$ (идеал $P+Q_j$ строго содержит
максимальный идеал $P$): найдутся $x_j\in P$, $y_j\in Q_j$ с $x_j+y_j=1$. Положим $e=y_1\cdots y_k$ (при $k=0$ полагаем $e=1$) и
$b=ae$. Тогда $e\in Q_j$ для всех $j$ и $e=\prod(1-x_j)\equiv1\pmod P$. Итак,
\[
b\equiv a\pmod P,\qquad b\in\sigma^{-1}(P)\ \text{ для всех }\sigma\in G\smallsetminus D_P .
\]
(Это частный случай китайской теоремы об остатках.)

\emph{Шаг 2: многочлен $g$.} Положим $g(X)=\prod_{\sigma\in G}(X-\sigma b)$. Его коэффициенты лежат в $A$ (\rf{4.27b}(a)) и
неподвижны относительно $G$: применение $\tau\in G$ лишь переставляет множители, так как $\sigma\mapsto\tau\sigma$~---
биекция $G$. Поэтому коэффициенты лежат в $A\cap\Q=\Z$ (\rf{4.27b}(b)), т.\,е.\ $g\in\Z[X]$ и $\bar g\in\F_p[X]$.

\emph{Шаг 3: редукция $g$.} Если $\sigma\notin D_P$, то $b\in\sigma^{-1}(P)$, т.\,е.\ $\sigma b\in P$ и
$\overline{\sigma b}=0$. Если $\sigma\in D_P$, то $\sigma b-\sigma a=\sigma(b-a)\in\sigma(P)=P$, т.\,е.\
$\overline{\sigma b}=\overline{\sigma a}=\bar\sigma(\bar a)$. Следовательно,
\[
\bar g(X)=X^{s}\prod_{\sigma\in D_P}\bigl(X-\bar\sigma(\bar a)\bigr),\qquad s=|G\smallsetminus D_P| .
\]

\emph{Шаг 4: подсчёт.} Множитель с $\sigma=\id$ показывает, что $\bar g(\bar a)=0$, поэтому минимальный многочлен $\mu$
элемента $\bar a$ над $\F_p$ делит $\bar g$. Пусть $\tau\in G_{\bar f}$. Тогда $\tau(\bar a)$~--- корень $\mu$, значит, и
$\bar g$; он ненулевой (так как $\bar a\ne0$), поэтому $\tau(\bar a)=\bar\sigma(\bar a)$ для некоторого $\sigma\in D_P$. Два
автоморфизма поля $\F_p[\bar a]$ над $\F_p$, совпадающие на $\bar a$, совпадают; значит, $\tau=\bar\sigma=\theta(\sigma)$.
(Здесь существенно, что $a\notin P$: иначе все корни $\mu$ могли бы <<спрятаться>> в множителе $X^s$.)

(b) \emph{Существование.} Группа $G_{\bar f}$ порождается автоморфизмом $x\mapsto x^p$ (\rf{4.25a}). По (a) существует
единственный $\sigma_P\in D_P$ с $\bar\sigma_P(x)=x^p$ для $x\in A/P$, т.\,е.\ $\overline{\sigma_P(a)}=\bar a^p$ для всех
$a\in A$; это и есть \eqref{eq:s6-frob}. Так как $\theta$~--- изоморфизм, а $G_{\bar f}$ порождается $\bar\sigma_P$, то
$D_P=\langle\sigma_P\rangle$.

\emph{Единственность во всей группе $G$.} Пусть $\tau\in G$ удовлетворяет \eqref{eq:s6-frob}. Если $a\in P$, то
$\tau(a)\equiv a^p\equiv0\pmod P$, т.\,е.\ $\tau(P)\subset P$; но $\tau(P)$~--- максимальный идеал, поэтому $\tau(P)=P$ и
$\tau\in D_P$. Далее $\bar\tau(\bar a)=\bar a^p=\bar\sigma_P(\bar a)$ для всех $a$, т.\,е.\ $\theta(\tau)=\theta(\sigma_P)$, и
$\tau=\sigma_P$ по инъективности $\theta$. (Короче: $\tau(\alpha_i)\equiv\alpha_i^p\equiv\sigma_P(\alpha_i)\pmod P$, оба
числа~--- корни $f$, а различные корни несравнимы по модулю $P$; значит, $\tau$ и $\sigma_P$ одинаково действуют на корнях.)
\end{proof}

\begin{keepfig}
\[
\begin{tikzcd}[column sep=large,row sep=large]
\Z\arrow[r,hook,cold]\arrow[d,two heads,cold] & A=\Z[\alpha_1,\dots,\alpha_m]\arrow[r,hook,cold]\arrow[d,two heads,cnew,"a\mapsto\bar a"] & E=\Q[\alpha_1,\dots,\alpha_m]\\
\F_p\arrow[r,hook,cnew] & A/P=\F_p[\bar\alpha_1,\dots,\bar\alpha_m] &
\end{tikzcd}
\]
\[
\begin{tikzcd}[column sep=large,row sep=large]
\{\alpha_1,\dots,\alpha_m\}\arrow[r,"\sigma",cold]\arrow[d,"r"',cnew] & \{\alpha_1,\dots,\alpha_m\}\arrow[d,"r",cnew]\\
\{\bar\alpha_1,\dots,\bar\alpha_m\}\arrow[r,"\bar\sigma"',cnew] & \{\bar\alpha_1,\dots,\bar\alpha_m\}
\end{tikzcd}
\]
\figcap{Вверху: кольцо $A$, порождённое корнями $f$, и его факторполе $A/P$~--- поле разложения $\bar f$ над $\F_p$ (\rf{4.27b},
\rf{4.27c}). Внизу: при $p\nmid D(f)$ редукция $r\colon\alpha_i\mapsto\bar\alpha_i$~--- биекция, и для $\sigma\in D_P$ квадрат
коммутативен: $r\circ\sigma=\bar\sigma\circ r$.}
\end{keepfig}

\begin{theorem}[4.28: теорема Дедекинда]
Пусть $f\in\Z[X]$~--- унитарный многочлен степени $m$, $p$~--- простое число, для которого $\bar f=f\bmod p$ не имеет
кратных корней \textup(равносильно, $p\nmid D(f)$, \rf{4.27c}\textup), и $\bar f=f_1\cdots f_r$, где $f_i\in\F_p[X]$
неприводимы степени $m_i$. Тогда $G_f$ содержит элемент $\sigma_f$, который как перестановка корней $f$ является
произведением $r$ независимых циклов длин $m_1,\dots,m_r$.
\end{theorem}

\begin{proof}
Пусть $P$~--- максимальный идеал $A$, содержащий $p$ (\rf{4.27b}(c)), и $\sigma_f=\sigma_P$~--- элемент Фробениуса
(\rf{4.27e}(b)). Редукция $r\colon\alpha_i\mapsto\bar\alpha_i$~--- биекция корней $f$ на корни $\bar f$ (\rf{4.27c}(c)), и для каждого
$\sigma\in D_P$
\[
r(\sigma\alpha_i)=\overline{\sigma\alpha_i}=\bar\sigma(\bar\alpha_i)=\bar\sigma(r(\alpha_i)),
\]
т.\,е.\ $\sigma=r^{-1}\circ\bar\sigma\circ r$ на корнях (квадрат на рисунке коммутативен). Сопряжённые перестановки имеют
одинаковый цикловой тип: $r$ переводит каждый цикл $(\beta_1\,\beta_2\,\dots\,\beta_k)$ перестановки $\sigma$ в цикл
$(r\beta_1\,r\beta_2\,\dots\,r\beta_k)$ перестановки $\bar\sigma$. Поэтому цикловой тип $\sigma_P$ на корнях $f$ равен
цикловому типу $\bar\sigma_P\colon x\mapsto x^p$ на корнях $\bar f$, а он по \rf{4.25a} равен $(m_1,\dots,m_r)$.
\end{proof}

Итак, утверждение <<действия $D_P$ на корнях $f$ и $G_{\bar f}$ на корнях $\bar f$ изоморфны>> почти тавтологично, но
опирается на два доказанных факта: редукция $r$ биективна (здесь нужно $p\nmid D(f)$), и $\theta\colon D_P\to G_{\bar f}$~---
изоморфизм. Если занумеровать корни $f$ и $\bar f$ согласованно ($\bar\alpha_i\leftrightarrow\alpha_i$), то $D_P$ и
$G_{\bar f}$~--- буквально одна и та же подгруппа $S_m$.

\begin{remark}[4.28a]
\begin{enumerate}[label=(\alph*),leftmargin=*]
\item Циклы элемента $\sigma_f$ соответствуют неприводимым множителям $\bar f$ над $\F_p$, а не множителям $f$ над $\Q$:
множители $f$ над $\Q$ отвечают орбитам \emph{всей} группы $G_f$ (\rf{4.27}), а циклы $\sigma_f$~--- орбитам лишь её подгруппы
$\langle\sigma_f\rangle=D_P$. Например, $X^5-X-1$ неприводим над $\Q$ (\rf{4.30}), так что $G_f$ транзитивна, но по модулю $2$
он раскладывается на множители степеней $2$ и $3$: у $\sigma_f$ две орбиты. Группа $G_f$, в отличие от
$G_{\bar f}\cong D_P$, вообще говоря, не циклическая; теорема Дедекинда утверждает лишь, что \emph{циклическая
подгруппа} $D_P$ в $G_f$ действует на корнях $f$ так же, как $G_{\bar f}$ на корнях $\bar f$. Крайний случай~---
многочлен $X^4-10X^2+1$ из \rf{1.18c}: он неприводим над $\Q$, но приводим по модулю каждого простого. Действительно, его
корни~--- $\pm\sqrt2\pm\sqrt3$, поле разложения~--- $\Q[\sqrt2,\sqrt3]$, и все неединичные элементы
$G_f\cong(\Z/2\Z)^2$ имеют порядок $2$; в $G_f$ нет $4$-циклов, и по (c) ниже $\bar f$ не может быть неприводим ни при
каком $p$.
\item Элемент $\sigma_f$ зависит от выбора $P$, но его цикловой тип~--- нет (он определяется разложением $\bar f$). О
зависимости от $P$~--- в \rf{4.28b}.
\item Порядок $\sigma_f$ равен наименьшему общему кратному чисел $m_1,\dots,m_r$; значит, это число делит $|G_f|$. Если
$\bar f$ неприводим при некотором $p$, то $f$ неприводим над $\Q$: разложение $f=gh$ в $\Z[X]$ (по лемме Гаусса, \rf{1.13}, \rf{1.14},
на унитарные множители) дало бы $\bar f=\bar g\bar h$. Теорема Дедекинда даёт больше: неприводимый над $\F_p$
многочлен сепарабелен ($\F_p$ совершенно), так что $p\nmid D(f)$ автоматически, и в $G_f$ есть цикл длины $m$. Если $\bar f$ раскладывается на линейные множители, то $\sigma_f=\id$ и $D_P=1$.
\end{enumerate}
\end{remark}

\subsection*{Элементы и классы Фробениуса}

Сколько максимальных идеалов $P\supset p$ в $A$ и как связаны их элементы Фробениуса?

\begin{proposition}[4.28b]
Пусть $p\nmid D(f)$.
\begin{enumerate}[label=\textup{(\alph*)}]
\item Группа $G$ транзитивно переставляет максимальные идеалы кольца $A$, содержащие $p$; поэтому их число равно
$(G:D_P)=|G|/[A/P:\F_p]$.
\item Для $\tau\in G$ имеем $\sigma_{\tau P}=\tau\sigma_P\tau^{-1}$.
\end{enumerate}
\end{proposition}

\begin{proof}
(a) Если $P\ni p$, то и $\tau P\ni\tau p=p$, и $\tau P$ максимален. Пусть $Q$~--- максимальный идеал, $p\in Q$, и
$Q\ne\sigma P$ для всех $\sigma\in G$. Пусть $P_1,\dots,P_k$~--- различные идеалы среди $\sigma P$; вместе с $Q$ это
различные максимальные идеалы. Конструкция шага 1 доказательства \rf{4.27e} даёт для каждого $l$ элемент $e_l\equiv1\pmod{P_l}$,
лежащий в $Q$ и во всех $P_{l'}$, $l'\ne l$. Тогда $x=e_1+\dots+e_k$ лежит в $Q$ и $x\equiv1\pmod{\sigma P}$ для всех
$\sigma\in G$. Число
$N=\prod_{\sigma\in G}\sigma(x)$ лежит в $A$ и неподвижно относительно $G$, т.\,е.\ $N\in A\cap\Q=\Z$; множитель с
$\sigma=\id$ лежит в $Q$, поэтому $N\in Q\cap\Z=p\Z\subset P$. Так как $P$ прост, некоторый множитель
$\sigma(x)\in P$, т.\,е.\ $x\in\sigma^{-1}P$; но $x\equiv1\pmod{\sigma^{-1}P}$, и $1\in\sigma^{-1}P$~--- противоречие. Итак,
все максимальные идеалы над $p$ образуют одну орбиту $G$, а стабилизатор $P$~--- это $D_P$, и $|D_P|=|G_{\bar f}|=[A/P:\F_p]$
по \rf{4.27e}.

(b) Для $a\in A$: $\tau\sigma_P\tau^{-1}(a)=\tau\bigl(\sigma_P(\tau^{-1}a)\bigr)$, а $\sigma_P(\tau^{-1}a)-(\tau^{-1}a)^p\in P$;
применяя $\tau$, получаем $\tau\sigma_P\tau^{-1}(a)-a^p\in\tau P$. По единственности в \rf{4.27e}(b) (для идеала $\tau P$)
$\tau\sigma_P\tau^{-1}=\sigma_{\tau P}$.
\end{proof}

\begin{definition}[4.28c: класс Фробениуса]
Пусть $p\nmid D(f)$. Множество $\{\sigma_P\mid P\supset p\}$ по \rf{4.28b}~--- это ровно один класс сопряжённости группы $G_f$; он
называется \emph{классом Фробениуса} в $p$ и обозначается $\Frob_p$. Все его элементы имеют цикловой тип
$(m_1,\dots,m_r)$ разложения $\bar f$. Если $G_f$ абелева, то $\Frob_p$ состоит из одного элемента $\sigma_p$, который
тогда называют элементом Фробениуса в $p$.
\end{definition}

Абелев случай разобран для круговых полей: в $\Gal(\Q(\zeta_n)/\Q)\cong(\Z/n\Z)^\times$ элемент Фробениуса в $p\nmid n$ есть
$\zeta\mapsto\zeta^p$, т.\,е.\ класс $[p]$ (\rf{5.10}, \rf{5.11j}). Это даёт и сам изоморфизм (второе доказательство \rf{5.10}), и
квадратичный закон взаимности (\rf{5.11l})~--- пример того, как классы Фробениуса применяются.

\begin{example}[4.28d: $X^3-2$]
Пусть $f=X^3-2$, $G_f\cong S_3$ (поле разложения $\Q[\sqrt[3]2,\zeta_3]$ имеет степень $6$), $D(f)=-108=-2^2\cdot3^3$.
Для простых $p\ge5$:
\begin{itemize}[leftmargin=*]
\item $p\equiv2\pmod3$: тогда $3\nmid p-1$, и $x\mapsto x^3$~--- автоморфизм циклической группы $\F_p^\times$ порядка $p-1$;
значит, у $\bar f$ ровно один корень в $\F_p$, и $\bar f=(\text{лин.})(\text{неприв. квадр.})$. Например,
$X^3-2\equiv(X-3)(X^2+3X+4)\pmod5$. Класс Фробениуса~--- три транспозиции; максимальных идеалов над $p$ три
($6/2$), и у каждого свой $\sigma_P$.
\item $p\equiv1\pmod3$: в $\F_p$ есть первообразный кубический корень из $1$, так что вместе с одним корнем $\bar f$
содержит все три. Либо $2$~--- куб в $\F_p$, и $\bar f$ раскладывается на линейные множители ($\Frob_p=\{\id\}$, шесть
идеалов; например, $p=31$: корни $4,7,20$), либо $\bar f$ неприводим ($\Frob_p$~--- два $3$-цикла, два идеала; например,
$p=7,13$).
\end{itemize}
Среди 2260 простых $5\le p<20000$ типы $(2,1)$, $(3)$, $(1,1,1)$ встречаются $1136$, $756$ и $368$ раз, т.\,е.\ с частотами
$\approx0{,}503$, $0{,}335$, $0{,}163$~--- близко к долям $3/6$, $2/6$, $1/6$ транспозиций, $3$-циклов и единицы в $S_3$
(ср.\ \rf{4.29a}). (Классический результат Гаусса: при $p\equiv1\pmod3$ число $2$~--- куб по модулю $p$ тогда и только тогда,
когда $p=x^2+27y^2$; так, $31=2^2+27$, $43=4^2+27$.)
\end{example}

\begin{aside}[4.29: кольцо целых и классы Фробениуса]
Пусть $E$~--- конечное расширение Галуа поля $\Q$ с группой $G$, и $\mathcal O_E$~--- \emph{кольцо целых} поля $E$:
множество элементов $E$, являющихся корнями унитарных многочленов из $\Z[X]$. Почему это кольцо? Элемент $\alpha$ цел
тогда и только тогда, когда абелева группа $\Z[\alpha]$ конечно порождена. (Если $\alpha$~--- корень унитарного
многочлена степени $n$, то $\Z[\alpha]$ порождается $1,\alpha,\dots,\alpha^{n-1}$, как в \rf{4.27b}(a). Обратно, если $\alpha$
лежит в подкольце $B\subset E$, конечно порождённом как абелева группа образующими $b_1,\dots,b_k$, то
$\alpha b_i=\sum_jc_{ij}b_j$, $c_{ij}\in\Z$; матрица $\alpha I-(c_{ij})$ переводит ненулевой вектор $(b_1,\dots,b_k)$ в ноль,
поэтому её определитель~--- унитарный многочлен от $\alpha$ с целыми коэффициентами~--- равен нулю.) Если $\alpha,\beta$
целы, то $B=\Z[\alpha,\beta]$ порождается конечным числом одночленов $\alpha^i\beta^j$, и $\alpha\pm\beta$, $\alpha\beta\in B$
целы по доказанному. Например, кольцо $A$ из \rf{4.27b} содержится в $\mathcal O_K$ для $K=\Q_f$, но может быть меньше его.

Пусть $P$~--- максимальный идеал $\mathcal O_E$ с $P\cap\Z=p\Z$. Тем же рассуждением, что в \rf{4.27e} и \rf{4.28b}, доказывается:
$G$ транзитивно переставляет такие идеалы; если $p$ \emph{не разветвлено} в $E$ (для $E=\Q_f$ это так, например, при
$p\nmid D(f)$), то существует единственный $\sigma_P\in G$ с $\sigma_PP=P$ и $\sigma_P(a)\equiv a^p\pmod P$ для всех
$a\in\mathcal O_E$~--- \emph{автоморфизм Фробениуса} в $P$; $\sigma_{\tau P}=\tau\sigma_P\tau^{-1}$, и класс сопряжённости
$\sigma_P$~--- класс Фробениуса в $p$. Для абелевой $G$ он состоит из одного элемента.

\emph{Геометрическая картина.} Максимальные идеалы кольца $\Z$~--- это $p\Z$, и их удобно представлять себе как точки
<<прямой>> $\operatorname{Spec}\Z$; максимальные идеалы $\mathcal O_E$~--- точки <<кривой>> $\operatorname{Spec}\mathcal O_E$,
а сопоставление $P\mapsto P\cap\Z$~--- отображение $\operatorname{Spec}\mathcal O_E\to\operatorname{Spec}\Z$, подобное
разветвлённому накрытию степени $|G|$. Слой над точкой $p$~--- идеалы $P\supset p$; группа $G$ действует на кривой и
транзитивно переставляет точки каждого слоя, как группа автоморфизмов накрытия Галуа. Над неразветвлённым $p$ в слое
$|G|/|D_P|$ точек, и каждая <<весит>> $|D_P|=[\mathcal O_E/P:\F_p]$; над разветвлённым точки <<склеиваются>> (так
$X^4+1\equiv(X+1)^4\pmod 2$). Группа разложения $D_P$~--- стабилизатор точки, а элемент Фробениуса~--- аналог
монодромии. Точка $p$ <<устроена как окружность>>: группа $\Gal(\Fpb/\F_p)$ порождена (топологически) Фробениусом
$x\mapsto x^p$ (\rf{7.17}), как фундаментальная группа окружности порождена одной петлёй; слой над $p$ накрывает эту
<<окружность>>, и обход петли переставляет его листы. Если отметить точку $P$ слоя, получится элемент $\sigma_P$ группы
$G$; при другом выборе точки~--- сопряжённый элемент. Строго эти аналогии оформляются в алгебраической теории чисел и
алгебраической геометрии.
\end{aside}

\begin{keepfig}
\centering
\begin{tikzpicture}[x=1cm,y=1cm,every node/.style={font=\small}]
\draw[cold,thick] (0,0) -- (10,0) node[right] {$\operatorname{Spec}\Z$};
\foreach \x/\l in {1/2,3/3,5/5,7/7,9/p} \fill[cold] (\x,0) circle (1.6pt) node[below=2pt] {$\l$};
\foreach \y in {1.2,1.7,2.2,2.7,3.2,3.7}{
  \draw[cnew,thin] (0,\y) .. controls (0.6,\y) and (0.75,2.45) .. (1,2.45);
  \draw[cnew,thin] (1,2.45) .. controls (1.25,2.45) and (1.4,\y) .. (2,\y) -- (10,\y);
}
\node[cnew,above right] at (10,3.7) {$\operatorname{Spec}\mathcal O_E$};
\fill[cnew] (1,2.45) circle (2.2pt);
\node[above] at (1,3.8) {разветвлено};
\foreach \y in {1.2,1.7,2.2,2.7,3.2,3.7} \fill[cnew] (5,\y) circle (1.6pt);
\node[above] at (5,3.8) {$D_P=1$};
\foreach \y in {1.45,2.45,3.45} \draw[cnew,thick] (9,\y) ellipse (0.18 and 0.38);
\node[above] at (9,3.85) {$|D_P|=2$};
\node[right] at (9.2,2.45) {$P$};
\draw[->,>=stealth] (5,1.05) -- (5,0.15);
\draw[->,>=stealth] (9,1.0) -- (9,0.15);
\end{tikzpicture}
\figcap{Отступление \rf{4.29}: <<накрытие>> $\operatorname{Spec}\mathcal O_E\to\operatorname{Spec}\Z$ степени $|G|=6$ (условная
картинка; это не поле разложения $X^3-2$ из \rf{4.28d}: там разветвлены $2$ и $3$, над $5$ лежат три точки с $|D_P|=2$, а
шесть точек~--- например, над $31$). Над разветвлённым простым (здесь $2$) листы склеиваются. Над неразветвлённым $p$
слой состоит из $|G|/|D_P|$ точек, и группа $G$ переставляет их транзитивно: над $5$~--- шесть точек ($D_P=1$), над
$p$~--- три <<толстые>> точки с полем вычетов $\mathcal O_E/P$ степени $|D_P|=2$ над $\F_p$. Элементы Фробениуса точек
одного слоя сопряжены.}
\end{keepfig}

\begin{aside}[4.29a: теорема Чеботарёва о плотности \textup(без доказательства\textup)]
Следующая глубокая теорема теории чисел здесь не доказывается; она используется только в пояснениях и в примерах
\rf{4.28d}, \rf{4.29b}, \rf{4.29c} и в алгоритме \rf{4.32b}, а доказательства остальных утверждений на неё не опираются.

Пусть $f\in\Z[X]$ унитарен и сепарабелен, $G=G_f$, и $C\subset G$~--- класс сопряжённости. Тогда множество простых
$p\nmid D(f)$ с $\Frob_p=C$ имеет плотность $|C|/|G|$:
\[
\lim_{x\to\infty}\frac{\#\{p\le x\mid p\nmid D(f),\ \Frob_p=C\}}{\#\{p\le x\}}=\frac{|C|}{|G|}.
\]
(Н.\,Г.~Чеботарёв, 1920-е годы; то же верно для любого конечного расширения Галуа $E/\Q$ и неразветвлённых $p$.) В
частности, каждый элемент $G$ является элементом Фробениуса $\sigma_P$ для бесконечно многих $p$. Для
$E=\Q(\zeta_n)$ это теорема Дирихле о простых числах в арифметических прогрессиях (её частный случай~--- \rf{5.11k}).

Для вычисления групп Галуа нужна только более слабая теорема плотности Фробениуса (1896): \emph{доля простых $p$, по
модулю которых $\bar f$ имеет данный тип разложения $(m_1,\dots,m_r)$, равна доле элементов $G_f$ циклового типа
$(m_1,\dots,m_r)$.} Иначе говоря, типы разложения $f$ по модулю простых <<распределены как цикловые типы случайного
элемента $G_f$>>. Существует и \emph{эффективный} вариант: если некоторый цикловой тип встречается в $G_f$, то он
встречается как тип разложения по модулю простого $p$, не превосходящего явной границы, зависящей от $f$ (при
обобщённой гипотезе Римана граница имеет порядок квадрата логарифма дискриминанта поля разложения).
\end{aside}

\begin{example}[4.29b]
\emph{Неприводимый унитарный многочлен $f\in\Z[X]$ степени $n\ge2$ не имеет корней по модулю бесконечного множества простых
чисел} (и даже множества положительной плотности). Действительно, при $p\nmid D(f)$ многочлен $\bar f$ не имеет корня
в $\F_p$ тогда и только тогда, когда у $\sigma_f$ нет циклов длины $1$ (\rf{4.28}), т.\,е.\ неподвижных точек. По теореме
Жордана в транзитивной группе $G\subset S_n$, $n\ge2$, есть перестановка без неподвижных точек: по лемме о числе орбит
среднее по $G$ число неподвижных точек элемента равно числу орбит, т.\,е.\ $1$, а у единицы их $n>1$; значит, у
какого-то элемента их меньше $1$, т.\,е.\ нет совсем. Если таких элементов в $G_f$ ровно $d\ge1$, то по \rf{4.29a} доля
простых, по модулю которых у $f$ нет корней, равна $d/|G_f|>0$. Например, $X^2+1$ не имеет корней по модулю простых
$p\equiv3\pmod4$ (доля $1/2$), а $X^3-2$~--- по модулю тех $p\equiv1\pmod3$, для которых $2$ не куб (доля $2/6$, \rf{4.28d}).
Без теоремы Чеботарёва бесконечность такого множества простых не очевидна.
\end{example}

\subsection*{Как угадать группу Галуа}

По теореме Дедекинда каждый тип разложения $\bar f$ ($p\nmid D(f)$) есть цикловой тип элемента $G_f$, а по теореме
плотности (\rf{4.29a}) типы встречаются с частотами, равными долям элементов соответствующего типа. Это позволяет быстро
<<угадать>> $G_f$; строгое доказательство того, что группа не больше, требует других средств (дискриминант, резольвенты).

\begin{example}[4.29c]
Транзитивные подгруппы $S_5$ (с точностью до сопряжения)~--- это $C_5$, $D_5$, $F_{20}$, $A_5$, $S_5$ (\rf{4.18a}); числа их
элементов каждого циклового типа сведены в таблицу после доказательства \rf{4.18a} (там же обозначения: $2\,1^3$~---
транспозиция, $2^2\,1$~--- произведение двух независимых транспозиций и т.\,д.).
Разложим по модулю всех простых $p<10\,000$, не делящих дискриминант (их $1227$), два многочлена: $f=X^5-X-1$, $D(f)=2869=19\cdot151$, и
$g=X^5-5X+12$, $D(g)=2^{12}5^6$. Наблюдаемые числа простых каждого типа и ожидаемые (доля в группе, умноженная на 1227):
\begin{center}
\small
\begin{tabular}{lccccccc}
\toprule
 & $1^5$ & $2\,1^3$ & $2^2\,1$ & $3\,1^2$ & $3\,2$ & $4\,1$ & $5$\\
\midrule
$f$ & 6 & 108 & 151 & 201 & 211 & 296 & 254\\
$S_5$: ожидается & 10{,}2 & 102{,}3 & 153{,}4 & 204{,}5 & 204{,}5 & 306{,}8 & 245{,}4\\
\midrule
$g$ & 116 & 0 & 625 & 0 & 0 & 0 & 486\\
$D_5$: ожидается & 122{,}7 & & 613{,}5 & & & & 490{,}8\\
\bottomrule
\end{tabular}
\end{center}
Для $f$ уже $p=2$ даёт тип $3\,2$, и этого достаточно для строгого вывода $G_f=S_5$ (\rf{4.30}). Для $g$ встречаются только типы
$1^5$, $2^2\,1$, $5$, в пропорции, близкой к $1:5:4$; это убедительно указывает на $G_g=D_5$, но не доказывает этого.
Строго из наблюдений следует лишь, что $G_g$ содержит $5$-цикл и элемент типа $2^2\,1$. Так как $D(g)=(2^65^3)^2$~---
квадрат, $G_g\subset A_5$ (\rf{4.2}); подгруппа $A_5$, содержащая элементы порядков $5$ и $2$, имеет порядок, делящийся на
$10$, а в $A_5$ нет подгрупп порядков $20$ и $30$ (индексов $3$ и $2$; см.\ доказательство \rf{4.18a}), так что $G_g$~--- это $D_5$ или $A_5$. Отсутствие в $G_g$ $3$-циклов
никакой конечный перебор простых не докажет; исключить $A_5$ позволяет резольвента Кэли; это сделано в
\rf{4.18g}(1) (другое рассуждение~--- в \rf{2.4d}).
\end{example}

\subsection*{Примеры с группой \texorpdfstring{$S_n$}{Sn}}

\begin{example}[4.30]
Пусть $f=X^5-X-1$. По формуле \rf{4.39} $D(X^5+aX+b)=5^5b^4+4^4a^5$, так что $D(f)=3125-256=2869=19\cdot151$; он не делится
на $2$ и $3$, и теорема Дедекинда применима при $p=2,3$.

\emph{По модулю $2$:} $X^5-X-1\equiv X^5+X+1=(X^2+X+1)(X^3+X^2+1)$ (раскройте скобки: $X^5+X^4+X^3+X^4+X^3+X^2+X^2+X+1$, и
в $\F_2$ удвоенные члены сокращаются). Множители неприводимы: многочлены степени $2$ и $3$ без корней в $\F_2$ (их
значения в $0$ и $1$ равны $1$) неприводимы, так как разложение содержало бы линейный множитель.

\emph{По модулю $3$:} $X^5-X-1\equiv X^5+2X+2$, и этот многочлен неприводим над $\F_3$. Корней в $\F_3$ нет: $a^3=a$, и
значит $a^5=a^3a^2=a^3=a$ для $a\in\F_3$, так что значение равно $a+2a+2=2\ne0$. Если бы он был приводим, то, не имея
линейных множителей, он разлагался бы в произведение неприводимых множителей степеней $2$ и $3$. Пусть $\alpha$~--- корень
квадратного множителя; тогда $\alpha\in\F_9$ (поле $\F_3[X]/(\text{множитель})$), так что $\alpha^9=\alpha$, и
$\alpha^5=-2\alpha-2=\alpha+1$. Отсюда
\[
\alpha=\alpha^9=\alpha^5\alpha^4=(\alpha+1)\alpha^4=\alpha^5+\alpha^4=\alpha^4+\alpha+1,
\]
т.\,е.\ $\alpha^4=-1$, и $\alpha^5=-\alpha$. Сравнивая с $\alpha^5=\alpha+1$, получаем $2\alpha=-1=2$, $\alpha=1\in\F_3$~---
но корней в $\F_3$ нет. Противоречие.

По теореме Дедекинда $G_f$ содержит перестановку типа $(2,3)$, т.\,е.\ $(i\,k)(l\,m\,n)$, и $5$-цикл. Независимые циклы
коммутируют, поэтому
\[
\bigl((i\,k)(l\,m\,n)\bigr)^3=(i\,k)^3(l\,m\,n)^3=(i\,k)\cdot\id=(i\,k).
\]
Группа $G_f\subset S_5$ содержит транспозицию и $5$-цикл, и $5$ простое; по \rf{4.15} $G_f=S_5$. (Неприводимость $f$ над $\Q$
следует уже из неприводимости по модулю $3$.)
\end{example}

\begin{lemma}[4.31]
Транзитивная подгруппа $H\subset S_n$, содержащая транспозицию и $(n-1)$-цикл, совпадает с $S_n$.
\end{lemma}

\begin{proof}
Перенумеровав точки, можно считать, что $(n-1)$-цикл~--- это $\rho=(1\,2\,\dots\,n{-}1)$; он оставляет на месте $n$. Будем
пользоваться правилом сопряжения: для перестановки $\pi$
\[
\pi\,(a\,b)\,\pi^{-1}=(\pi(a)\;\pi(b)),
\]
так как $\pi(a)\xmapsto{\pi^{-1}}a\xmapsto{(a\,b)}b\xmapsto{\pi}\pi(b)$, симметрично $\pi(b)\mapsto\pi(a)$, а точка, отличная
от $\pi(a),\pi(b)$, переходит сама в себя.

\emph{Шаг 1: транспозиция вида $(i\,n)$.} Пусть $(a\,b)\in H$~--- данная транспозиция. По транзитивности есть $\pi\in H$ с
$\pi(b)=n$. Тогда $\pi(a\,b)\pi^{-1}=(\pi(a)\;n)\in H$, и $i=\pi(a)\ne n$, так как $\pi$ инъективна и $a\ne b$. (Если
$n\in\{a,b\}$, то шаг не нужен.)

\emph{Шаг 2: все $(j\,n)$.} Так как $\rho^k(n)=n$, то $\rho^k(i\,n)\rho^{-k}=(\rho^k(i)\;n)$, а $\rho^k(i)$ при $k=0,\dots,n-2$
пробегает все числа $1,\dots,n-1$. Значит, $(1\,n),(2\,n),\dots,(n-1\;n)\in H$.

\emph{Шаг 3.} Эти транспозиции порождают $S_n$: $(j\,k)=(j\,n)(k\,n)(j\,n)$ (по правилу сопряжения с $\pi=(j\,n)$), а
транспозиции порождают $S_n$. Итак, $H=S_n$.
\end{proof}

\begin{example}[4.32: многочлены с группой $S_n$]
Следующий рецепт даёт многочлены степени $n\ge3$ над $\Q$ с группой Галуа $S_n$. Выберем унитарные многочлены
$f_1,f_2,f_3\in\Z[X]$ степени $n$, сепарабельные по модулю соответствующих простых, с такими разложениями:
\begin{enumerate}[label=(\alph*)]
\item $f_1$ неприводим по модулю $2$;
\item по модулю $3$ многочлен $f_2$ равен (многочлен степени $1$)$\cdot$(неприводимый многочлен степени $n-1$);
\item по модулю $5$ многочлен $f_3$ равен (неприводимый многочлен степени $2$)$\cdot$(произведение одного или двух
неприводимых многочленов нечётной степени).
\end{enumerate}
Положим
\[
f=-15f_1+10f_2+6f_3 .
\]
Тогда $f$ унитарен ($-15+10+6=1$), и
\[
f\equiv f_1\pmod2,\qquad f\equiv f_2\pmod3,\qquad f\equiv f_3\pmod5,
\]
так как $-15\equiv1$, $10\equiv6\equiv0\pmod2$; $10\equiv1$, $-15\equiv6\equiv0\pmod3$; $6\equiv1$, $-15\equiv10\equiv0\pmod5$.
Поэтому $\bar f$ сепарабелен по
модулю $2$, $3$ и $5$, т.\,е.\ $2,3,5\nmid D(f)$, и теорема Дедекинда применима:
\begin{enumerate}[label=(\roman*)]
\item $G_f$ содержит $n$-цикл (из (a)); значит, $G_f$ транзитивна (и $f$ неприводим);
\item $G_f$ содержит $(n-1)$-цикл (из (b));
\item $G_f$ содержит перестановку $\tau\sigma$, где $\tau$~--- транспозиция, а $\sigma$~--- произведение одного или двух
независимых от $\tau$ циклов нечётных длин (из (c)). Пусть $k$~--- порядок $\sigma$ (наименьшее общее кратное нечётных
длин); $k$ нечётно, $\tau$ и $\sigma$ коммутируют, и $(\tau\sigma)^k=\tau^k\sigma^k=\tau$. Итак, $G_f$ содержит транспозицию.
\end{enumerate}
По лемме \rf{4.31} $G_f=S_n$.

\emph{Почему рецепт выполним.} Над каждым $\F_p$ есть неприводимые многочлены любой степени (\rf{4.22c}(b)), так что
нужные разложения существуют. В (c) при нечётном $n-2$ берём один неприводимый множитель степени $n-2$, а при чётном
$n-2\ge2$~--- два: степеней $1$ и $n-3$ (если $n-3=1$, т.\,е.\ $n=4$, это два \emph{различных} линейных множителя, скажем
$X$ и $X-1$, чтобы $\bar f_3$ был сепарабелен). Взаимная простота различных неприводимых множителей обеспечивает
сепарабельность. В качестве $f_i$ берутся любые унитарные целочисленные поднятия этих многочленов.

Главная идея: каждое простое число даёт только <<нижнюю оценку>> группы (элемент в ней), и простые подбираются так,
чтобы найденных элементов хватило для порождения $S_n$; по лемме \rf{4.31} достаточно трёх типов: $n$, $(n-1)+1$ и
$2+(\text{нечётные})$. Коэффициенты $-15,10,6$~--- решение китайской задачи об остатках для модулей $2,3,5$ (они
сравнимы с $1$ по <<своему>> модулю и с $0$ по двум другим), а условие $-15+10+6=1$ делает $f$ унитарным.
\end{example}

\begin{example}[4.32a]
Проведём построение при $n=6$. Возьмём
\[
f_1=X^6+X+1,\qquad f_2=X(X^5+2X+1),\qquad f_3=X(X^2+2)(X^3+X+1).
\]
\begin{itemize}[leftmargin=*]
\item $X^6+X+1$ неприводим над $\F_2$. Корней нет (значения в $0,1$ равны $1$). Единственный неприводимый
квадратный многочлен над $\F_2$~--- $X^2+X+1$; его корень $\omega$ удовлетворяет $\omega^3=1$, и
$\omega^6+\omega+1=\omega\ne0$. Неприводимый кубический множитель имел бы корень $\beta\in\F_8$, $\beta^7=1$, и тогда
$0=\beta^6+\beta+1=\beta^{-1}+\beta+1$, т.\,е.\ $\beta^2+\beta+1=0$ и $\beta\in\F_4$; но $\F_4\cap\F_8=\F_2$ (\rf{4.21}), а
в $\F_2$ корней нет.
\item $X^5+2X+1=X^5-X+1$ неприводим над $\F_3$~--- так же, как в \rf{4.30}: значения в $\F_3$ равны $a-a+1=1$; если $\alpha\in\F_9$~---
корень, то $\alpha^5=\alpha-1$, $\alpha=\alpha^9=(\alpha-1)\alpha^4=\alpha^5-\alpha^4=\alpha-1-\alpha^4$, откуда $\alpha^4=-1$,
$\alpha^5=-\alpha=\alpha-1$, $2\alpha=1$, $\alpha=2\in\F_3$~--- противоречие. Множитель $X$ взаимно прост с ним.
\item По модулю $5$: $X^2+2$ неприводим ($-2\equiv3$~--- не квадрат: квадраты в $\F_5$ суть $0,1,4$), $X^3+X+1$
неприводим (его значения в $0,1,2,3,4$ равны $1,3,1,1,4$), и множители $X$, $X^2+2$, $X^3+X+1$ попарно различны.
Это тип $2+(1+3)$, как в (c).
\end{itemize}
Получаем
\[
f=-15f_1+10f_2+6f_3=X^6+18X^4+6X^3+32X^2+7X-15,
\]
и $G_f=S_6$. (Для контроля: $D(f)=383\cdot487\cdot3989\cdot216577$ не делится на $2,3,5$, и разложения $f$ по модулю $2,3,5$
имеют типы $6$, $5+1$, $3+2+1$.)
\end{example}

\subsection*{Общая стратегия}

\begin{remark}[4.32b: алгоритм вычисления $G_f$]
Пусть $f\in\Q[X]$; умножив на подходящее число и заменив $X$ на $X/c$, можно считать, что $f\in\Z[X]$ унитарен (корни
умножаются на $c$, поле разложения и группа не меняются).
\begin{enumerate}[leftmargin=*]
\item Разложить $f$ на неприводимые множители над $\Q$ (\rf{1.18} и далее); группы приводимых многочленов сводятся к
группам множителей (\rf{3.27e}). Далее $f$ неприводим степени $n$, и $G_f$~--- транзитивная подгруппа $S_n$, порядок
которой делится на $n$ (\rf{3.27d}).
\item Вычислить $D(f)$ (например, через результант, \rf{4.36}) и проверить, квадрат ли он: $G_f\subset A_n$ тогда и только
тогда, когда $D(f)$~--- квадрат (\rf{4.2}).
\item Разложить $f$ по модулю последовательных простых $p\nmid D(f)$ и записать типы разложения~--- цикловые типы
элементов $G_f$ (\rf{4.28}). По \rf{4.29a} каждый тип, встречающийся в $G_f$, рано или поздно появится (и в эффективном
варианте~--- для $p$ меньше явной границы), с частотой, равной доле таких элементов.
\item Сравнить с таблицей транзитивных подгрупп $S_n$ (порядка, делящегося на $n$) и их цикловых типов. Если
найденные элементы порождают одну из групп (как в \rf{4.30}, \rf{4.32}), группа найдена строго. Если же наблюдаемые типы
указывают на меньшую группу, доказывать, что больше ничего нет, нужно иначе: резольвенты (\rf{4.9}--\rf{4.13} для $n=4$; \rf{4.18c}--\rf{4.18g}
для $n=5$), действие $G_f$ на множествах $r$-элементных подмножеств корней
(разложение соответствующей резольвенты степени $\binom nr$), дискриминант.
\end{enumerate}
Цикловых типов и длин орбит на $r$-подмножествах почти всегда достаточно, чтобы различить транзитивные подгруппы $S_n$
для небольших $n$; их таблицы составлены для $n\le11$ (Батлер и Маккей, 1983; см.\ также Сойчер и Маккей, 1985). Системы
компьютерной алгебры реализуют эту программу; например, в PARI/GP команда \texttt{polgalois(f)} для неприводимого $f$
степени $\le11$ выдаёт порядок группы, знак ($+1$, если $G_f\subset A_n$), номер и название группы. Так,
для \texttt{x\^{}5-x-1} она выдаёт порядок $120$, знак $-1$ и название \texttt{S5}, а для \texttt{x\^{}5-5*x+12}~--- порядок
$10$, знак $+1$ и название \texttt{D(5) = 5:2}.
\end{remark}

\begin{aside}[4.33]
Для унитарного многочлена степени $n$ с ограниченными целыми коэффициентами ожидается, что его группа Галуа равна $S_n$
с вероятностью, стремящейся к $1$ при $n\to\infty$. Доказана (Л.~Бари-Сорокер и Г.~Козма, 2020) более слабая версия:
если коэффициенты $a_0,\dots,a_{n-1}$ многочлена $X^n+a_{n-1}X^{n-1}+\dots+a_0$ выбираются независимо и равномерно из
множества $\{1,\dots,L\}$, где $L$ делится хотя бы на четыре различных простых числа (например, $L=210$), то с
вероятностью, стремящейся к $1$, многочлен неприводим и его группа Галуа содержит $A_n$. Отличить $A_n$ от $S_n$
(т.\,е.\ доказать, что дискриминант почти никогда не является квадратом) этот результат не позволяет.
\end{aside}

\begin{aside}[4.34]
Из $(2H+1)^n$ унитарных многочленов $X^n+a_1X^{n-1}+\dots+a_n\in\Z[X]$ с $\max|a_i|\le H$ сколько имеют группу Галуа, отличную
от $S_n$? Их заведомо $\gg H^{n-1}$: таковы, например, все многочлены с $a_n=0$ (они приводимы). Ван дер Варден
предположил (1936), а Бхаргава доказал (2021), что их число есть $O(H^{n-1})$, т.\,е.\ эта оценка точна по порядку. Таким
образом, $S_n$~--- группа Галуа <<почти всех>> многочленов, и все остальные группы встречаются как исключения; это ещё
одно объяснение, почему так легко построить примеры с группой $S_n$ (\rf{4.17}, \rf{4.30}, \rf{4.32}).
\end{aside}

\section{Приложение: вычисление дискриминантов с помощью результантов}

Пусть $f,g\in F[X]$ и в некотором поле разложения $fg$
\[
f(X)=a\prod_{i=1}^n(X-\alpha_i),\qquad g(X)=b\prod_{j=1}^m(X-\beta_j),\qquad ab\ne0 .
\]
\emph{Результант} $f$ и $g$:
\[
\operatorname{Res}(f,g)=a^mb^n\prod_{i,j}(\alpha_i-\beta_j).
\]
Он равен нулю тогда и только тогда, когда у $f$ и $g$ есть общий корень, т.\,е.\ $\gcd(f,g)\ne1$. Для ненулевой константы
$g=b$ ($m=0$) имеем $\operatorname{Res}(f,b)=b^n$. Хотя определение использует корни, результант лежит в $F$: по \rf{4.35}(b)
ниже он равен $a^m\prod_ig(\alpha_i)$, а это значение симметрического многочлена $\prod_ig(X_i)\in F[X_1,\dots,X_n]$ в
корнях унитарного многочлена $a^{-1}f\in F[X]$, т.\,е.\ элемент $F$ по \rf{5.37}(c). (С точностью до знака результант равен
определителю матрицы Сильвестра из \rf{1.8f}; здесь это не используется.)

\begin{proposition}[4.35]
В этих обозначениях:
\begin{enumerate}[label=\textup{(\alph*)}]
\item $\operatorname{Res}(f,g)=(-1)^{mn}\operatorname{Res}(g,f)$;
\item $\operatorname{Res}(f,g)=a^m\prod_{i=1}^ng(\alpha_i)$;
\item если $g\equiv g_1\pmod f$ в $F[X]$, где $g_1\ne0$ и $\deg g_1=m_1$, то
$\operatorname{Res}(f,g)=a^{m-m_1}\operatorname{Res}(f,g_1)$; если же $n\ge1$ и $f\mid g$, то $\operatorname{Res}(f,g)=0$.
\end{enumerate}
\end{proposition}

\begin{proof}
(a) $\operatorname{Res}(g,f)=b^na^m\prod_{i,j}(\beta_j-\alpha_i)$, и каждый из $mn$ множителей отличается знаком от
множителя в $\operatorname{Res}(f,g)$.
(b) $g(\alpha_i)=b\prod_j(\alpha_i-\beta_j)$, так что $a^m\prod_ig(\alpha_i)=a^mb^n\prod_{i,j}(\alpha_i-\beta_j)$.
(c) $g=g_1+hf$, поэтому $g(\alpha_i)=g_1(\alpha_i)$, и по (b), применённому к $g$ и к $g_1$,
\[
\operatorname{Res}(f,g)=a^m\prod_ig(\alpha_i)=a^{m-m_1}\cdot a^{m_1}\prod_ig_1(\alpha_i)=a^{m-m_1}\operatorname{Res}(f,g_1).
\]
Если $f\mid g$, то $g(\alpha_i)=0$ для всех $i$, и при $n\ge1$ по (b) $\operatorname{Res}(f,g)=0$.
\end{proof}

Эти формулы позволяют вычислять результанты, как алгоритм Евклида вычисляет наибольший общий делитель: (c) понижает
степень второго многочлена (берём $g_1$~--- остаток от деления $g$ на $f$), (a) меняет многочлены местами, и так до
тех пор, пока один из многочленов не станет линейным или константой; тогда ответ даёт (b):
\[
\operatorname{Res}\bigl(c(X-\gamma),g\bigr)=c^{\deg g}g(\gamma),\qquad \operatorname{Res}(f,b)=b^{\deg f},
\]
а нулевой остаток означает, что $\operatorname{Res}=0$.

\begin{proposition}[4.36]
Пусть $f\in F[X]$~--- унитарный многочлен степени $n$ с $f'\ne0$. Тогда
\[
D(f)=(-1)^{\frac{n(n-1)}2}\operatorname{Res}(f,f').
\]
Если $\deg f'=n-1$ \textup(т.\,е.\ $\operatorname{char}F\nmid n$; например, $\operatorname{char}F=0$\textup), то также
$D(f)=(-1)^{\frac{n(n-1)}2}\operatorname{Res}(f',f)$.
\end{proposition}

\begin{proof}
Пусть $f=\prod_{i=1}^n(X-\alpha_i)$. По \rf{4.35}(b) с $a=1$ имеем $\operatorname{Res}(f,f')=\prod_if'(\alpha_i)$, и первое
равенство~--- это \rf{4.2d}(a). Если $\deg f'=n-1$, то по \rf{4.35}(a)
$\operatorname{Res}(f',f)=(-1)^{n(n-1)}\operatorname{Res}(f,f')=\operatorname{Res}(f,f')$.
\end{proof}

(Если $f'=0$, то результант не определён, а $D(f)=0$: каждый корень $f$~--- корень $f'$, т.\,е.\ кратный.)

\emph{Характеристика, делящая $n$.} Пусть $\operatorname{char}F=p\mid n$ и $f'\ne0$. Тогда $m'=\deg f'<n-1$, и по \rf{4.35}(a)
$\operatorname{Res}(f',f)=(-1)^{nm'}\operatorname{Res}(f,f')$; при нечётном $nm'$, $\operatorname{char}F\ne2$ и $D(f)\ne0$
второе равенство \rf{4.36} \emph{неверно}. Пример: $f=X^3+X^2+2$ над $\F_3$ (неприводим: $f(0)=2$, $f(1)=1$, $f(2)=2$). Здесь
$f'=2X$, $m'=1$, $\alpha_1\alpha_2\alpha_3=-2=1$, и
\[
\operatorname{Res}(f,f')=\prod_i2\alpha_i=8=2,\qquad \operatorname{Res}(f',f)=2^3f(0)=16=1
\]
(второе~--- по \rf{4.35}(b) для $f'=2(X-0)$). Поэтому $D(f)=-\operatorname{Res}(f,f')=-2=1$, что согласуется с формулой \rf{4.2a}
($D=-4\cdot2-27\cdot4=-116\equiv1\pmod3$), а $-\operatorname{Res}(f',f)=-1=2\ne D(f)$. Кстати, $D(f)=1$~--- квадрат, как и
должно быть (\rf{4.2}(b)): группа Галуа неприводимого кубического многочлена над конечным полем циклическая порядка $3$
(\rf{4.25a}), т.\,е.\ $G_f=A_3$.

Если $\operatorname{char}F\nmid n$ и $f'=n\prod_{j=1}^{n-1}(X-c_j)$, то по \rf{4.35}(b), применённому к $f'$ (старший
коэффициент $n$), $\operatorname{Res}(f',f)=n^n\prod_jf(c_j)$, и \rf{4.36} превращается в формулу критических значений
\rf{4.2d}(b). Для неприводимого $f$ произведение $\prod_if'(\alpha_i)$ есть норма $\operatorname{Nm}_{F[\alpha]/F}(f'(\alpha))$ (\rf{5.49}). Для
неунитарного $f=a_0\prod_{i=1}^n(X-\alpha_i)$ с $\deg f'=n-1$ то же вычисление даёт
$\operatorname{Res}(f,f')=a_0^{\,n-1}\prod_if'(\alpha_i)=(-1)^{\frac{n(n-1)}2}a_0\cdot a_0^{2n-2}\prod_{i<j}(\alpha_i-\alpha_j)^2$
(ср.\ \rf{4.2b}(d)).

\emph{О делениях в примерах.} Ниже мы делим на коэффициенты и на целые числа ($3$, $2b$, $n$, $(n-1)a$ и т.\,п.), т.\,е.\
фактически считаем, что $\operatorname{char}F=0$ и эти величины ненулевые. Получающиеся формулы верны над любым полем и
при любых значениях коэффициентов по принципу специализации \rf{4.2c}.

\begin{example}[4.37]
Пусть $f=X^3+bX+c$, $b\ne0$. По \rf{4.36} ($n=3$, $(-1)^3=-1$) $D(f)=-\operatorname{Res}(3X^2+b,\,X^3+bX+c)$. Деление с остатком
\[
X^3+bX+c=\tfrac X3(3X^2+b)+\tfrac{2b}3X+c
\]
и \rf{4.35}(c) (старший коэффициент $3$, $m-m_1=3-1=2$) дают
\[
D(f)=-3^2\operatorname{Res}\bigl(3X^2+b,\ \tfrac{2b}3X+c\bigr)=-3^2\operatorname{Res}\bigl(\tfrac{2b}3\bigl(X+\tfrac{3c}{2b}\bigr),\ 3X^2+b\bigr)
\]
(\rf{4.35}(a), $(-1)^{2}=1$). По \rf{4.35}(b), подставляя корень $-\frac{3c}{2b}$:
\[
D(f)=-3^2\Bigl(\frac{2b}3\Bigr)^2\Bigl(3\Bigl(\frac{3c}{2b}\Bigr)^2+b\Bigr)=-4b^2\Bigl(\frac{27c^2}{4b^2}+b\Bigr)=-4b^3-27c^2 .
\]
При $b=0$ формула тоже верна: $D(X^3+c)=-27c^2$ (\rf{4.2c}). Другие выводы той же формулы собраны в \rf{4.2e}; вывод
\rf{4.2e}(iv) через произведения $x_if'(x_i)$~--- по существу то же вычисление, что и здесь, только без слова
<<результант>>.
\end{example}

\begin{example}[4.38]
Пусть $f=X^5+X+1$. По \rf{4.36} ($n=5$, $(-1)^{10}=1$) $D(f)=\operatorname{Res}(5X^4+1,\,X^5+X+1)$. Так как
$X^5+X+1=\frac X5(5X^4+1)+\frac45X+1$, по \rf{4.35}(c) ($m-m_1=5-1=4$) и \rf{4.35}(a) ($(-1)^4=1$)
\begin{align*}
D(f)&=5^4\operatorname{Res}\bigl(5X^4+1,\ \tfrac45X+1\bigr)=5^4\operatorname{Res}\bigl(\tfrac45\bigl(X+\tfrac54\bigr),\ 5X^4+1\bigr)\\
&=5^4\Bigl(\frac45\Bigr)^4\Bigl(5\Bigl(\frac54\Bigr)^4+1\Bigr)=5^5+4^4=3381 .
\end{align*}
Здесь $3381=3\cdot7^2\cdot23$. Заметим, что многочлен $X^5+X+1=(X^2+X+1)(X^3-X^2+1)$ приводим (его корни~--- в том числе
первообразные кубические корни $\omega$ из $1$: $\omega^5+\omega+1=\omega^2+\omega+1=0$), так что это пример лишь на
вычисление дискриминанта. Для приводимого многочлена дискриминант можно найти и по множителям. Если $g,h$ унитарны с
корнями $\alpha_i$ и $\beta_j$, то пары корней произведения $gh$ бывают трёх видов (оба~--- корни $g$, оба~--- корни $h$,
по одному), поэтому
\[
D(gh)=D(g)\,D(h)\prod_{i,j}(\alpha_i-\beta_j)^2=D(g)\,D(h)\operatorname{Res}(g,h)^2 .
\]
Проверка: $D(X^2+X+1)=-3$, $D(X^3-X^2+1)=-23$ (\rf{4.2a}), $\operatorname{Res}(X^2+X+1,\,X^3-X^2+1)=(2-\omega^2)(2-\omega)=7$
(\rf{4.35}(b); $\omega,\omega^2$~--- корни $X^2+X+1$, и $\omega^3-\omega^2+1=2-\omega^2$), и $(-3)(-23)\cdot7^2=3381$.

Поучительно сравнить $X^5+X+1$ с $X^5-X-1$ из \rf{4.30}: по модулю $2$ эти многочлены совпадают и
равны $(X^2+X+1)(X^3+X^2+1)$, но первый приводим над $\Q$, а второй неприводим. Разложение по модулю $p$ даёт
элемент группы Галуа, но не разложение над $\Q$.
\end{example}

\begin{example}[4.39]
Пусть $f=X^n+aX+b$, $n\ge2$. Тогда
\[
D(f)=(-1)^{\frac{n(n-1)}2}\Bigl(n^nb^{n-1}+(-1)^{n-1}(n-1)^{n-1}a^n\Bigr).
\]
Вычисление по той же схеме (при $a\ne0$, в поле, где $n(n-1)\ne0$): $X^n+aX+b=\frac Xn(nX^{n-1}+a)+\frac{n-1}naX+b$, и
\begin{align*}
D(f)&=(-1)^{\frac{n(n-1)}2}\operatorname{Res}(nX^{n-1}+a,\ X^n+aX+b)\\
&=(-1)^{\frac{n(n-1)}2}n^{n-1}\operatorname{Res}\Bigl(nX^{n-1}+a,\ \tfrac{n-1}na\Bigl(X+\tfrac{nb}{(n-1)a}\Bigr)\Bigr)
\qquad\text{(4.35(c))}\\
&=(-1)^{\frac{n(n-1)}2}(-1)^{n-1}n^{n-1}\Bigl(\frac{(n-1)a}n\Bigr)^{n-1}\Bigl(n\Bigl(\frac{-nb}{(n-1)a}\Bigr)^{n-1}+a\Bigr)
\qquad\text{(4.35(a), (b))}\\
&=(-1)^{\frac{n(n-1)}2}(-1)^{n-1}\Bigl(n^n(-b)^{n-1}+(n-1)^{n-1}a^n\Bigr)\\
&=(-1)^{\frac{n(n-1)}2}\Bigl(n^nb^{n-1}+(-1)^{n-1}(n-1)^{n-1}a^n\Bigr).
\end{align*}
(Множитель $n^{n-1}$ в первой строке~--- это $(\text{старший коэффициент } f')^{\,n-1}$ из \rf{4.35}(c), а $(-1)^{n-1}$~--- знак из
\rf{4.35}(a) при степенях $n-1$ и $1$.) По \rf{4.2c} формула верна всегда. Частные случаи: $n=2$: $a^2-4b$; $n=3$: $-4a^3-27b^2$ (как в \rf{4.37}); $n=4$:
$256b^3-27a^4$ (\rf{4.2f}(b)); $n=5$: $D(X^5+aX+b)=5^5b^4+4^4a^5$, в частности $D(X^5+X+1)=3381$ (как в \rf{4.38}) и
$D(X^5-X-1)=3125-256=2869$ (\rf{4.30}). Ту же формулу дают критические значения \rf{4.2d}(b) (критические точки удовлетворяют
$c^{n-1}=-\frac an$, и $f(c)=\frac{n-1}n\,ac+b$~--- то же вычисление) и норма (\rf{5.50}).
\end{example}
